% Modernised reading — fonds Alexandre Grothendieck, shelfmark 29, batch 10
% (pages 181-200).
%
% This is an INTERPRETATION, not a transcription. It re-reads the pages in
% today's notation and vocabulary, states the hypotheses the manuscript leaves
% implicit, and drops the critical apparatus so that the mathematics reads
% straight through. Everything the brackets and underlines carried moves into a
% footnote. In any dispute of substance the transcription governs: whoever
% cites Grothendieck must cite batch-10.fr.tex.
%
% Scope: the folder was read WHOLE before any of this was written — all eleven
% batches, pages 1-216, in one pass — and one file was then written per batch.
% This is the batch where that scope pays for itself. Pages 190-200 answer, in
% the same carton, the question page 24 asked in 1967: how to normalise the
% notion of ramification datum so that the auxiliary object disappears. The
% answer — a champ over the étale site with a continuous geometric realisation
% — is exactly what page 24 predicted one would be « infailliblement conduit »
% to. And the three examples of 1967 reappear, in order, as three CONSEQUENCES:
% Prop. 3 and its corollary give example 1, page 195 gives example 2, and page
% 196's remark that R is not generated by its global sections is why example 3
% had to exist at all. This reading names each junction where it occurs.
%
% NOTATION, and the one real trap of this folder. The 1967 datum is written
% calligraphic R = (C, s) throughout the earlier files; the object of these
% pages is written roman (R, r). They are DIFFERENT: the first is auxiliary and
% the second is the champ itself. Grothendieck reuses the letter and calls both
% « donnée » or « domaine » de ramification; the distinction is stated here at
% the point where the two meet.
%
% Departures from the page, each footnoted where it occurs:
%   - the definition of « subordonné » on page 192 is damaged; it is fixed by
%     page 194's own use of it, and the reading says so;
%   - page 186's proposition carries an illegible word in its first hypothesis
%     and the statement as legible does not visibly give quasi-compactness; it
%     is recorded, not repaired;
%   - R4 on page 190 is only partly legible and is given as far as it reads.
%
% Pass: Opus 5 (claude-opus-5), 2026-08-29 — first pass, unchecked against the
% pages by a human.
\documentclass[11pt,a4paper]{article}
% Le préambule commun à toutes les transcriptions du fonds Grothendieck.
%
% Il tient en une idée : **l'incertitude se compose, elle ne se résout pas.**
% Une transcription qui lisse un mot douteux a détruit la seule chose pour
% laquelle on la faisait — un lecteur doit pouvoir savoir, à chaque endroit,
% ce qui était sur le feuillet et ce qui a été ajouté. Les six macros qui
% suivent sont donc l'appareil critique, et non de la décoration.
%
% Les mêmes commandes sont reconnues par `scripts/render.mjs`, qui produit la
% vue de lecture du site. Une macro ajoutée ici doit l'être aussi là-bas,
% faute de quoi le rendu échouera bruyamment — ce qui est voulu : un rendu
% silencieusement faux est pire qu'un rendu absent.

\ProvidesPackage{grothendieck}[2026/08/08 Transcriptions du fonds Grothendieck]

\usepackage{amsmath,amssymb,amsthm}
\usepackage{mathrsfs}
\usepackage{stmaryrd}
% Les diagrammes commutatifs sont partout dans ce fonds : c'est la langue
% même de la géométrie algébrique de Grothendieck, et une transcription qui
% les rendrait en prose ne transcrirait rien.
\usepackage{tikz-cd}
% Français par défaut — c'est la langue des feuillets — mais l'anglais est
% chargé aussi : la traduction et le résumé sont des documents anglais, et la
% typographie française (espaces avant les deux-points) y serait une faute.
% Ces documents basculent par \selectlanguage{english} en tête de corps.
\usepackage[english,french]{babel}
\usepackage{fontspec}
\usepackage{geometry}
\geometry{a4paper,margin=25mm,marginparwidth=28mm}
\usepackage{marginnote}
\usepackage{xcolor}
% ulem, not soul. `soul`'s \st and \ul are built on a character-by-character
% scan that XeLaTeX's font machinery breaks: the document fails with an
% undefined control sequence rather than degrading. ulem does the same two jobs
% and is Unicode-engine safe. `normalem` stops it hijacking \emph, which the
% transcriptions use for Grothendieck's own emphasis.
\usepackage[normalem]{ulem}
\usepackage{hyperref}
\hypersetup{colorlinks=true,linkcolor=black,urlcolor=black,pdfborder={0 0 0}}

% --- Métadonnées du lot -----------------------------------------------------
% Reprises telles quelles par le script de rendu, qui les affiche en tête de
% la vue de lecture. Elles fixent ce que le fichier prétend couvrir : sans
% elles, une transcription flotte sans point d'attache dans un fonds de seize
% mille feuillets.

\newcommand{\folder}[1]{\gdef\grothFolder{#1}}
\newcommand{\batch}[1]{\gdef\grothBatch{#1}}
\newcommand{\leaves}[2]{\gdef\grothFirst{#1}\gdef\grothLast{#2}}
\newcommand{\foldertitle}[1]{\gdef\grothFoldertitle{#1}}
\gdef\grothFolder{?}\gdef\grothBatch{?}\gdef\grothFirst{?}\gdef\grothLast{?}\gdef\grothFoldertitle{}

% --- Le résumé --------------------------------------------------------------
%
% Il ouvre la lecture modernisée, sous le titre « Résumé », et il est composé
% en petit. Ce n'est pas un ornement : c'est le seul passage du document qui
% ne soit pas des mathématiques, et un lecteur doit pouvoir le reconnaître
% avant de l'avoir lu — puis le sauter s'il connaît déjà le sujet, ou s'y
% arrêter s'il ne le connaît pas. Le filet qui le suit marque la reprise du
% corps.

\newenvironment{resume}
  {\small\setlength{\parskip}{0.45em}}
  {\par\medskip\hrule\medskip}

% --- Mots-clés ---------------------------------------------------------------
%
% En anglais, en clôture du résumé de la lecture modernisée : le vocabulaire
% moderne sous lequel chercher ce que les feuillets construisent. Le manifeste
% du site (`npm run manifest`) les extrait pour étiqueter la cote entière —
% cette ligne est la source unique des tags, il n'y a pas de fichier à côté.
\newcommand{\keywords}[1]{\par\smallskip\noindent
  {\footnotesize\textsc{Keywords} --- \textit{#1}}\par}

% --- L'appareil critique ----------------------------------------------------

% Le numéro de feuillet, en marge. C'est l'ancre qui permet de régler un
% désaccord sur une formule en regardant *un* feuillet plutôt qu'une chemise
% de six cents. Le site s'en sert aussi pour tourner les pages du fac-similé
% quand on fait défiler la transcription.
\newcommand{\leaf}[1]{%
  \marginnote{\footnotesize\color{gray!70}\textsf{#1}}%
  \label{leaf:#1}\ignorespaces}

% Illisible : on ne devine pas. Un mot inventé qui se lit comme les autres est
% le pire résultat possible.
\newcommand{\ill}{\textcolor{red!60!black}{[\,\ldots\,]}}

% Lecture douteuse : le mot est donné, et signalé comme incertain.
\newcommand{\uncertain}[1]{\uline{#1}}

% Ajout de l'éditeur — une lettre manquante, un mot sous-entendu.
\newcommand{\add}[1]{\textcolor{blue!50!black}{[#1]}}

% Biffé par Grothendieck lui-même. Ce qu'il a rayé fait partie du document :
% une rature dit souvent où la pensée a changé de direction.
\newcommand{\struck}[1]{\sout{#1}}

% Note du transcripteur, sur l'état matériel du feuillet.
\newcommand{\note}[1]{\par\smallskip{\small\color{gray!60!black}\textit{[#1]}}\par\smallskip}

% Note marginale de Grothendieck, distincte du corps du texte.
\newcommand{\marginal}[1]{\par\smallskip\noindent\hspace{1em}%
  {\small\color{gray!60!black}#1}\par\smallskip}

% --- Titre ------------------------------------------------------------------

\AtBeginDocument{%
  \begin{center}
    {\large\bfseries Fonds Alexandre Grothendieck --- cote n\textsuperscript{o}~\grothFolder}\\[2pt]
    {\itshape\grothFoldertitle}\\[4pt]
    {\small feuillets \grothFirst--\grothLast\ (lot \grothBatch)}\\[2pt]
    {\footnotesize\color{gray!60!black}Transcription établie d'après le fac-similé numérisé
     par l'Université de Montpellier.\\
     Les lectures incertaines sont soulignées, les illisibles notées [\,\ldots\,],
     les ajouts entre crochets.}
  \end{center}
  \vspace{1em}}

\endinput


\folder{29}
\batch{10}
\pages{181}{200}
\dating{1959-1969}
\watermark{Édition de démonstration\\interprétation personnelle de l'œuvre}
\foldertitle{Groupe fondamental [Autour de SGA 4 et SGA 7] — lecture modernisée}

\begin{document}

\section*{Résumé}

\begin{resume}
Une chemise de carton brun porte, souligné deux fois, « Domaines de
ramification / Catégories multigaloisiennes ». Ce qu'elle contient est la
réponse à une question que ce carton posait à sa page 24, deux cents pages plus
tôt et en 1967.

Voici la question. Pour définir la catégorie des revêtements admettant un type
de ramification donné, on avait dû se donner un catalogue de modèles ---
Grothendieck l'appelait une \emph{donnée de ramification} --- et il avait
aussitôt reconnu que ce catalogue était arbitraire : deux catalogues très
différents pouvaient définir la même catégorie. Il avait alors écrit, comme on
note une difficulté qu'on remet à plus tard : si l'on veut éliminer
l'arbitraire, « il semble qu'on soit infailliblement conduit à prendre comme
catégorie la catégorie des revêtements elle-même qu'on se propose de
construire ». Et il avait donné l'analogie juste : le catalogue est au
résultat ce qu'un \emph{site} est à son \emph{topos}.

Ces pages font exactement cela. Un \emph{domaine de ramification} sur un schéma
$S$ n'est plus un catalogue : c'est un \emph{champ} au-dessus du site étale de
$S$ --- c'est-à-dire une famille de catégories, une par ouvert étale, qui se
recollent --- dont chaque fibre est une catégorie multigaloisienne, munie d'une
« réalisation géométrique » vers les schémas. Quatre axiomes, R1 à R4, disent
ce que la réalisation géométrique doit respecter, et le premier travail est de
montrer qu'ils font des familles surjectives une \emph{topologie}. L'objet
auxiliaire a disparu : il ne reste que la catégorie, et le foncteur qui la
projette sur la géométrie.

Ce qui suit est le plus remarquable, et c'est ce qu'une lecture du carton
entier permet de voir : \emph{les trois exemples de 1967 reviennent, dans
l'ordre, comme trois conséquences}. La proposition 3 et son corollaire
retrouvent l'exemple 1 --- un schéma muni d'un groupe fini d'opérateurs ---
comme la description du sous-champ engendré par un objet galoisien. La page 195
retrouve l'exemple 2 --- un système projectif de tels objets --- comme la
description de $R$ tout entier quand $S$ est affine. Et la page 196 explique
enfin pourquoi l'exemple 3, si laborieux, était indispensable : le champ $R$
n'est pas en général engendré par ses sections globales, de sorte qu'il faut
localiser et \emph{recoller}. Un procédé de recollement est donné, puis une
construction générale par les \emph{gerbes}.

Le lot se termine sur les exemples kummériens, et il faut les regarder. La
gerbe des racines $n_{i}$-ièmes des équations d'une famille de diviseurs
$D_{i}$ a une obstruction dans $H^{2}$, et Grothendieck la nomme : c'est « la
classe de cohomologie des $D_{i}$ ». C'est le théorème des pages 89 et 100,
énoncé une quatrième fois. Et la gerbe elle-même, avec les revêtements
équivariants de son revêtement kummérien, est ce qu'on appelle aujourd'hui le
\emph{champ des racines} : le groupe fondamental modéré d'un couple $(S, D)$
est le groupe fondamental d'un champ.

Les noms modernes à chercher sont \emph{champ}, \emph{gerbe}, \emph{champ des
racines}, \emph{2-catégorie}, \emph{géométrie logarithmique}.

\keywords{ramification domain, gerbe, root stack, 2-category, stack of topoi}
\end{resume}

\pagerange{181}{185}
\section*{Foncteurs semi-cosimpliciaux, et le van Kampen (pages 181, 182, 184 et 185)}

Ces quatre pages continuent le calcul de la page 180 en le rendant abstrait :
au lieu du diagramme simplicial d'un recouvrement, un diagramme de foncteurs.
Elles sont d'un crayon très pâle sur un feuillet dont le verso transparaît, et
la transcription y laisse beaucoup d'illisible.

\subsection*{Quand un foncteur est pro-représentable}

Soit $F$ un endofoncteur de la catégorie des ensembles. On cherche quand
\[
  F(I) \;\simeq\; C(K, I) ,
\]
$C(K,I)$ désignant les applications continues d'un espace $K$ totalement
discontinu dans l'ensemble discret $I$ --- c'est-à-dire les fonctions
localement constantes à valeurs dans $I$. Trois conditions nécessaires :

\begin{enumerate}
  \item[a)] $F$ commute aux limites projectives finies, c'est-à-dire est exact
    à gauche ;
  \item[b)] $F\bigl(\bigcap_{\alpha} I_{\alpha}\bigr) =
    \bigcap_{\alpha} F(I_{\alpha})$ pour toute famille filtrante de
    sous-ensembles de $I$ ;
  \item[c)] $\varprojlim_{j} \pi_{0j} \to \pi_{0i}$ est \emph{surjectif} pour
    tout $i$.
\end{enumerate}

Et le point : \emph{ces conditions sont nécessaires et suffisantes}, et
équivalent à dire que $F$ est \emph{strictement pro-représentable} par un
système projectif strict d'ensembles $(\pi_{0i})$ ; on prend alors
$K = \varprojlim \pi_{0i}$ muni de la topologie limite des topologies discrètes.

La condition c) est celle qui distingue un système projectif \emph{strict} d'un
système quelconque : elle dit qu'aucune information n'est perdue en passant à
la limite. Sans elle, $K$ pourrait être vide sans qu'aucun $\pi_{0i}$ le soit.

Une remarque en note, avec sa référence : si $F$ est représenté de deux façons
par des $C(L, I)$ avec $L$ un groupe topologique, on a bien
$\pi_{0}(L) \to K$, mais si $L$ n'est \emph{pas compact} ce n'est ni un
isomorphisme, ni surjectif, ni injectif.\footnote{La référence portée en marge
est « Bourb.\ Topologie II, p.\ 232, prop.\ 6 ». C'est la mise en garde qui
justifie tout le soin apporté à la stricte pro-représentabilité : le groupe
topologique qui « représente » n'est déterminé que sous une hypothèse de
compacité, et c'est pourquoi le carton parle de \emph{pro}-objets et non de
groupes.}

\subsection*{Le noyau, et le foncteur semi-cosimplicial}

Soit $F^{*} : F^{0} \rightrightarrows F^{1}$ un couple de foncteurs, et
\[
  \pi^{0}(F^{*}) \;=\; \operatorname{Ker}\bigl(F^{0} \rightrightarrows
  F^{1}\bigr) .
\]
Si $F^{0}$ et $F^{1}$ sont pro-représentables par $K_{0}$ et $K_{1}$, les deux
flèches proviennent de $K_{0} \leftleftarrows K_{1}$, et $\pi^{0}(F^{*})$ l'est
aussi.

On considère alors un \emph{foncteur semi-cosimplicial}

\begin{tikzcd}[column sep=small, row sep=small]
F^0 \arrow[r, bend left=12] \arrow[r, bend right=12]
  & F^1 \arrow[r, bend left=20] \arrow[r] \arrow[r, bend right=20] & F^2
\end{tikzcd}

auquel on associe deux foncteurs $\pi^{0}(F^{*}; E)$ et $\pi^{1}(F^{*}; G)$,
pour $E$ un ensemble et $G$ un groupe variables ; et, si l'on se donne de plus
un \emph{pointage} $\xi : F^{*} \to f^{*}$,

\begin{tikzcd}[column sep=small, row sep=small, nodes={font=\scriptsize}]
F^0 \arrow[r, bend left=12] \arrow[r, bend right=12] \arrow[d, "\xi^0"]
  & F^1 \arrow[r, bend left=20] \arrow[r] \arrow[r, bend right=20]
    \arrow[d, "\xi^1"] & F^2 \arrow[d, "\xi^2"] \\
f^0 \arrow[r, bend left=12] \arrow[r, bend right=12]
  & f^1 \arrow[r, bend left=20] \arrow[r] \arrow[r, bend right=20] & f^2
\end{tikzcd}

les versions pointées $\pi^{0}(F^{*}, \xi; E)$ et $\pi^{1}(F^{*}, \xi; G)$.
Lorsque tout est pro-représentable, on se ramène par un procédé standard à un
\emph{ensemble semi-simplicial pointé}
\[
  K_{0} \leftleftarrows K_{1} \Lleftarrow K_{2} ,
  \qquad
  k_{0} \leftleftarrows k_{1} \Lleftarrow k_{2} ,
\]
et l'on retrouve exactement les cocycles de la page 180. C'est le $\pi^{0}$ et
le $\pi^{1}$ non abéliens d'un objet semi-cosimplicial.

\subsection*{Le van Kampen}

La page 185 se termine sur le programme, et il faut le lire :

\begin{quote}
Il y a lieu de généraliser ces constructions --- on a vu que des $K_{i}$, des
$\overline{K}_{i}$, on a des catégories galoisiennes à chaque point, pour
décider de descente : le \emph{Van Kampen} \dots
\end{quote}

C'est-à-dire : la descente d'un revêtement le long d'un recouvrement est un
calcul de $\pi^{0}$ et de $\pi^{1}$ sur le nerf de ce recouvrement, et le
théorème de van Kampen --- qui reconstruit le groupe fondamental d'un espace à
partir de ceux d'un recouvrement --- en est la forme classique. Le carton s'en
servait déjà librement à la page 60 ; il en donne ici la raison.\footnote{Le
cadre général qu'appellent ces pages --- pro-objets, ensembles semi-simpliciaux,
$\pi^{i}$ non abéliens d'un objet cosimplicial --- est celui de l'homotopie
étale d'Artin et Mazur, dont le mémoire paraît en 1969. Le carton n'est pas
daté sur ces pages et cette lecture n'en tire pas de chronologie.}

\pagerange{183}{187}
\section*{Trois feuillets d'un autre ordre (pages 183, 186 et 187)}

\subsection*{Un feuillet d'esquisses}

La page 183 est un feuillet en travers, couvert de tracés sans ordre de
lecture : le diagramme $\overline{K}_{0} \leftleftarrows \overline{K}_{1}
\Lleftarrow \overline{K}_{2}$ avec les mentions « fini » et « groupe
topologique discret » ; la formule d'équivalence des données de descente de la
page 180 ; l'anneau $k[[s,t]]$ ; un chemin d'arêtes de $x$ à $a$ dont les
flèches ne vont pas toutes dans le même sens ; et, encadré, un symbole
$\pi_{1}^{\#}(S,a)$ surmonté de « $S$ connexe ». Le croisillon est tracé deux
fois de la même façon, donc c'est une notation ; rien sur la feuille n'en fixe
le sens, et cette lecture ne le conjecture pas.

\subsection*{Une proposition sur les morphismes de type fini}

La page 186 est écrite au net, à l'encre bleue, et porte une proposition sans
rapport avec ce qui l'entoure.

\textbf{Proposition.} Soit $f : X \to Y$ un morphisme avec 1) $Y$ localement
noethérien ; 2) $X$ n'ayant qu'un nombre fini de composantes irréductibles ;
3) $f$ localement de type fini et \emph{séparé}. Alors $f$ est \emph{de type
fini}.

La démonstration se ramène à $Y$ affine noethérien et $X$ intègre, puis
construit un $X''$ fini au-dessus de $X$ et birationnel au-dessus d'un $X'$
intègre, de sorte que $X'' \to X$ soit fini et surjectif ; il suffit alors de
voir $X''$ de type fini sur $Y$, et $X'' \to X'$ étant birationnel entre
schémas intègres, on conclut par la normalisation.\footnote{L'hypothèse 1)
porte, entre « $Y$ loc.\ noeth. » et « de type fini », un mot que la
transcription ne lit pas, et les hypothèses 1) et 3) sont l'une et l'autre
biffées et réécrites sur la page. Telle qu'elle se laisse lire, la proposition
n'assure pas visiblement la quasi-compacité de $X$, qui est ce qu'il faut
démontrer ; c'est sans doute là que le mot manquant se trouve. La proposition
est donc enregistrée ici comme la page la porte, et non réparée : le seul
endroit d'où la quasi-compacité pourrait venir est l'hypothèse illisible.}

\subsection*{Un opérateur d'homotopie pour $d''$}

La page 187 est barrée de deux longues diagonales et se lit néanmoins d'un bout
à l'autre. Elle est d'analyse complexe --- le seul feuillet du carton qui le
soit --- et donne un opérateur d'homotopie explicite pour l'opérateur de
Cauchy--Riemann $d''$ en $n$ variables.

Sur un polydisque, on dispose pour chaque variable de deux opérateurs
intégraux : le noyau de Cauchy, qui projette sur les fonctions holomorphes en
$z_{j}$,
\[
  U_{j} f \;=\; \frac{1}{2\pi i} \int_{B_{j}}
    f(z_{1}, \ldots, t, \ldots, z_{n})\, \frac{dt}{t - z_{j}} ,
\]
et l'opérateur de Cauchy--Pompeiu, qui résout $\partial/\partial\bar{z}_{j}$,
\[
  T_{j} f \;=\; \frac{1}{2\pi i} \iint_{K_{j}}
    f(z_{1}, \ldots, t, \ldots, z_{n})\,
    \frac{dt \wedge \overline{dt}}{t - z_{j}} .
\]
On introduit ensuite, à partir des fonctions symétriques élémentaires
$p_{k} = \sum x_{1}x_{2}\cdots x_{k}$, la combinaison
\[
  \varphi_{n}(x_{1}, \ldots, x_{n}) \;=\; \frac{1}{n}\left(1 +
  \frac{p_{1}}{n-1} + \frac{p_{2}}{\binom{n-1}{2}} + \cdots + p_{n-1}\right) ,
\]
et l'on note $\varphi_{n}(U ; j_{1}, \ldots, j_{q})$ ce que devient
$\varphi_{n}(U_{1}, \ldots, U_{n})$ quand on remplace par zéro les
$U_{j_{1}}, \ldots, U_{j_{q}}$. L'opérateur d'homotopie, pour une forme
$\omega = a\, d\bar{z}_{j_{1}} \wedge \cdots \wedge d\bar{z}_{j_{q}}$ avec
$q \geqslant 1$, est alors
\[
  P\omega \;=\; \varphi_{n}(U ; j_{1}, \ldots, j_{q}) \ast
  \sum_{\gamma=1}^{q} (-1)^{\gamma-1} (T_{j_{\gamma}} a)\;
  d\bar{z}_{j_{1}} \wedge \cdots \wedge
  \widehat{d\bar{z}_{j_{\gamma}}} \wedge \cdots \wedge d\bar{z}_{j_{q}} ,
\]
et il vérifie
\[
  P\, d'' f \;=\; f - U_{1}U_{2}\cdots U_{n}\, f ,
\]
c'est-à-dire : \emph{$f$ moins sa partie holomorphe est un bord}. C'est le
lemme de Dolbeault--Grothendieck, avec son opérateur écrit
explicitement.\footnote{Le signe qui sépare $\varphi_{n}$ de la somme est un
astérisque appuyé que la page n'explique pas ; la transcription le signale et
cette lecture ne lui suppose pas de sens. La combinaison $\varphi_{n}$ est ce
qui rend l'opérateur indépendant de l'ordre choisi sur les indices.}

\pagerange{190}{191}
\section*{Domaines de ramification : les axiomes R1 à R4 (pages 190 et 191)}

Grothendieck pagine cette suite 1 à 10, puis 11 à 16 dans le lot suivant.

\subsection*{La définition}

Soit $S$ un schéma, $\mathrm{Et}(S)$ son site étale et $\mathcal{U}$ un
univers. Un \emph{domaine de ramification} sur $S$ est un couple $(R, r)$ formé
d'un \emph{champ en $\mathcal{U}$-catégories multigaloisiennes} $R$ sur
$\mathrm{Et}(S)$ et d'un morphisme \emph{continu}
\[
  r : R \longrightarrow \mathrm{Fl}_{S} ,
\]
appelé \emph{réalisation géométrique}, où $\mathrm{Fl}_{S}$ est la catégorie
fibrée sur $\mathrm{Et}(S)$ dont la fibre en $S'$ est $(\mathrm{Sch})_{/S'}$.
Un objet $X$ de $R_{S'}$ est un \emph{revêtement de type $R$ sur $S'$}, et
$r(X)$ est sa réalisation géométrique.

\begin{quote}
\textbf{Il faut s'arrêter ici.} Le carton emploie les mêmes mots pour deux
objets différents. La \emph{donnée} de ramification de 1967 --- notée
$\mathcal{R} = (\mathcal{C}, s)$ dans les lectures des lots 1 et 2 --- est un
catalogue auxiliaire : une catégorie fibrée quelconque, avec un foncteur vers
les schémas, dont on \emph{déduit} la catégorie des revêtements. Le
\emph{domaine} de ramification $(R, r)$ de cette page \emph{est} la catégorie
des revêtements, prise comme donnée, et c'est un champ. Les deux ne sont pas de
même nature, et le second est ce que la page 24 annonçait qu'on serait
« infailliblement conduit » à prendre.
\end{quote}

Les axiomes sont les suivants.

\begin{itemize}
\item[\textbf{R1}] $r$ commute aux \emph{sommes finies} et au \emph{passage au
  quotient} par une relation d'équivalence. Il n'est en général \emph{pas}
  exact à gauche.

\item[\textbf{R2}] Pour $S' \in \mathrm{Et}(S)$ et $X \in R_{S'}$,
  l'application $X' \mapsto r(X')$ est une \emph{bijection} de l'ensemble des
  sous-objets directs de $X$ sur l'ensemble des parties ouvertes et fermées de
  $r(X)$.

\item[\textbf{R3}] Pour une famille $(X_{i})$ --- pas nécessairement finie ---
  de sous-objets de $X$ telle que les $r(X_{i})$ recouvrent $r(X)$, la famille
  $X_{i} \to X$ est une famille de descente \emph{effective}.

\item[\textbf{R4}] Pour $X \to Y$ dominant dans $R_{S}$, on veut que les
  $r(X_{i})$ recouvrent $r(Y)$.\footnote{L'axiome R4 est le plus abîmé des
  quatre sur la page ; la transcription en laisse une partie illisible. Il est
  donné ici tel qu'il se lit, et non complété. Sa fonction se déduit néanmoins
  de l'usage qui en est fait : R1 et R3 servent à faire de la surjectivité une
  topologie, R2 à identifier sous-objets et parties ouvertes et fermées, et R4
  à assurer que la surjectivité se propage --- c'est ce que la page 191 note :
  « on a utilisé seulement R1 ; R3 signifie que les topologies sur $R_{S'}$
  sont localement finies ; R4 que \dots\ L'axiome sur le changement de base est
  trivial. »}
\end{itemize}

Les marges de la page portent trois desiderata qui ne sont pas des axiomes mais
disent ce qu'on voudrait : que les foncteurs de changement de base soient
\emph{exacts} ; que dans $\mathrm{Fl}_{S}$ le quotient soit déterminable ; et
que la famille des objets de $R_{S'}$ soit localement constante finie.

\subsection*{Proposition 1 : c'est une topologie}

Disons qu'une famille de flèches $X_{i} \to X$ de $R_{S'}$ est
\emph{couvrante} si la famille $r(X_{i}) \to r(X)$ est surjective ---
c'est-à-dire, grâce à R1 et « c'est Kif Kif », si les images $X_{i} \subset X$
sont telles que les $r(X_{i})$ recouvrent $r(X)$.

\textbf{Proposition 1.} Ces familles sont les familles couvrantes d'une
\emph{topologie} sur $R_{S'}$, et le foncteur de changement de base
$R_{S'} \to R_{S''}$ est \emph{continu}.

La vérification des quatre axiomes d'une topologie est page 191, et elle est
courte :

\begin{enumerate}
  \item[a)] une famille majorée par une famille couvrante est couvrante ---
    trivial ;
  \item[b)] une somme est couvrante --- trivial ;
  \item[c)] stabilité par changement de base : grâce à b) on se ramène au cas
    où les $X_{i} \to X$ sont des monomorphismes ; mais alors $r$ commute aux
    produits fibrés $X_{i} \times_{X} X'$, et l'on gagne ;
  \item[d)] nature locale : si des $X'_{\alpha} \to X$ couvrent $X$ et que pour
    chaque $\alpha$ les $X_{i} \times_{X} X'_{\alpha}$ couvrent
    $X'_{\alpha}$, alors les $X_{i}$ couvrent $X$ --- même réduction aux
    monomorphismes.
\end{enumerate}

Le mécanisme est le même dans c) et d), et il vaut d'être nommé : \emph{$r$
n'est pas exact à gauche en général, mais il l'est sur les monomorphismes},
et b) permet toujours de s'y ramener. C'est là tout le contenu de R1 et R2.

\textbf{Corollaire.} Prenant les topos associés aux $R_{S'}$, on trouve un
\emph{champ en topos} $\tilde{R}_{S'}$ au-dessus de la sous-catégorie des
objets localement constants finis de $R_{S'}$.

Un bloc rayé de plusieurs diagonales pose ensuite une question à laquelle la
page ne répond pas : pour $X \to Y$ avec $Y = \coprod_{n} Y_{n}$ et chaque
$X \times_{Y} Y_{n} \to Y_{n}$ de somme effective, la famille des
$Y_{n} \to Y$ est-elle couvrante ? Elle est laissée ouverte.

\pagerange{192}{194}
\section*{Objets subordonnés, et l'exemple 1 de 1967 retrouvé (pages 192 à 194)}

\subsection*{Vide et connexe}

\textbf{Proposition 2} (utilise R1 et R2). Pour $X \in R_{S'}$ :
\[
  X = \varnothing \;\Longleftrightarrow\; r(X) = \varnothing ,
  \qquad
  X \text{ connexe} \;\Longleftrightarrow\; r(X) \text{ connexe} .
\]

\subsection*{Subordination}

Soit $C$ une catégorie multigaloisienne et $P$ un objet. Un objet $X$ de $C$
est dit \emph{subordonné} à $P$ si $X_{P} \to P$ est \emph{constant} --- au
sens où $X_{P}$ est somme de copies de $P$.\footnote{La définition est abîmée
sur la page 192 et la transcription y laisse deux passages illisibles. Elle est
fixée ici par l'usage qu'en fait la page 194, où les objets subordonnés à $P$
sont mis en bijection avec les objets « constants par morceaux » sur $r(P)$ ;
c'est cette phrase-là qui décide, non le tracé de la page 192. La condition de
connexité que la page ajoute est celle sous laquelle « constant » et
« constant par morceaux » coïncident.}

Pour $P$ fixé, les objets subordonnés à $P$ forment une sous-catégorie
\emph{pleine} stable par limites projectives finies et par facteurs directs ---
car les objets constants par morceaux sur $P$ forment une telle sous-catégorie
de $C_{/P}$ --- donc une \emph{catégorie multigaloisienne pleine} $C^{(P)}$
de $C$.

Prenant $C = R_{S'}$ et les objets \emph{localement} subordonnés à $P$, on
obtient une sous-catégorie fibrée pleine multigaloisienne $R^{P}$ de $R$.

\subsection*{R2 équivaut à la fidélité de $r$}

Un feuillet inséré entre ses pages 3 et 4, et rayé de plusieurs diagonales,
porte une proposition qu'il vaut la peine de retenir :

\textbf{Proposition.} Sous R1, R3 et R4, l'axiome \textbf{R2 équivaut à ce que
$r$ soit \emph{fidèle}}.

\emph{Démonstration.} Supposons $r$ fidèle. Soient $X'$, $X''$ deux sous-objets
de $X$ avec $r(X') \supset r(X'')$ dans $r(X)$. Comme
$r(X' \cap X'') = r(X') \cap r(X'')$, les inclusions
$X' \cap X'' \subset X'$ et $X' \cap X'' \subset X''$ deviennent des
isomorphismes par $r$, ce qui ramène au cas d'une inclusion
$X' \subset X''$, soit $X'' = X' \sqcup X_{1}$. Alors
$r(X'') = r(X') \sqcup r(X_{1})$, et $r(X') = r(X'')$ équivaut à
$r(X_{1}) = \varnothing$, donc à $X_{1} = \varnothing$.

Cette proposition mérite d'être mise en regard de tout le lot 1. En 1967, la
\emph{non-}pleine-fidélité de la réalisation géométrique était le fait
fondamental --- « il importera de bien distinguer entre un revêtement de type
$\mathcal{R}$ et sa réalisation géométrique » --- et la \emph{fidélité} était
une propriété qu'on vérifiait au cas par cas sur les exemples. Ici, elle est
promue au rang d'axiome, sous la forme R2, et c'est le seul axiome des quatre
qui porte sur ce que $r$ ne perd pas.

\subsection*{Proposition 3, et l'exemple 1 de 1967}

\textbf{Proposition 3.} Soit $P$ un objet \emph{galoisien de groupe $G$} dans
$R_{S'}$. Le foncteur $X \mapsto r(X \times P)/r(P)$ est une \emph{équivalence}
de la catégorie des objets de $R_{S'}$ subordonnés à $P$ avec la catégorie des
objets constants par morceaux sur $r(P)$ munis d'une action de $G$ compatible
avec celle qu'il a sur $r(P)$.

La démonstration est en deux temps : utilisant R2, le foncteur est pleinement
fidèle ; utilisant R2 et R3, il est essentiellement surjectif. Et comme $G$
opère sur $R$ il opère sur $r(P)$, l'équivalence étant compatible.

\textbf{Corollaire 1.} Soit $P$ galoisien de groupe $G$ dans $R_{S}$. Le même
foncteur induit une équivalence de la catégorie des objets de $R_{S}$
\emph{localement} subordonnés à $P$ avec la catégorie des revêtements étales de
$Z = r(P)$ munis d'une action de $G$ compatible, satisfaisant la condition
suivante --- automatiquement vérifiée si $Z$ est fini sur $S$ : les $S' \to S$
étales tels que $Z_{S'}$ y devienne constant par morceaux recouvrent $S$.

Grothendieck ajoute lui-même, après un trait oblique : « Localisant sur $S$, on
trouve une description du champ $R^{P}$ en termes des schémas $Z$ à groupe
d'opérateurs $G$ sur $S$. »

\begin{quote}
\emph{C'est l'exemple 1 du tapuscrit de 1967}, page 14 du carton : un schéma
$Z$ sur $S$ muni d'un groupe fini $G$ d'opérateurs, la catégorie des
revêtements étales de $Z$ à opérateurs, et $r(Z') = Z'/G$. Ce qui était le
premier et le plus simple des trois exemples devient ici la description du
sous-champ engendré par un objet galoisien. L'exemple est devenu un théorème.
\end{quote}

\pagerange{195}{196}
\section*{L'exemple 2, et pourquoi l'exemple 3 devait exister (pages 195 et 196)}

\subsection*{Proposition 4 : la réciproque}

\textbf{Proposition 4.} Pour tout schéma $Z$ sur $S$ muni d'un groupe fini $G$
d'opérateurs agissant de façon \emph{admissible}, le champ décrit ci-dessus est
un champ en $\mathcal{U}$-catégories multigaloisiennes, et le foncteur
« passage au quotient par $G$ » --- qui est défini, l'action étant admissible
--- en fait un \emph{domaine de ramification} sur $S$.

La vérification est brève, et il faut la lire pour voir combien les axiomes
sont peu coûteux. Que $R_{S'}$ soit multigaloisienne : on part de la catégorie
des revêtements à opérateurs $G$ de $Z_{S'}$, dont on prend une sous-catégorie
pleine stable par limites projectives finies et par facteurs directs. Que R1
soit vérifié : $r$ est un \emph{passage au quotient}, donc commute aux limites
inductives finies, calculées au sens des schémas. R2, R3 et R4 sont
immédiats.\footnote{Le mot que la transcription lit « donnée » dans la
conclusion de la proposition 4 pourrait être « domaine » --- c'est le terme que
la page 190 définit et que la chemise porte en titre --- la lettre ne tranchant
pas et l'article féminin allant contre. C'est le seul endroit du carton où le
choix entre les deux mots a une conséquence, puisqu'il s'agit précisément de
dire que la construction produit l'objet nouveau et non l'objet auxiliaire de
1967. Le contexte impose « domaine » et c'est ce qui est écrit ici.}

\subsection*{L'exemple 2}

Supposons $S$ affine et $R$ soumis à « une petite condition supplémentaire »
vérifiée en pratique. Alors tout objet de $R_{S}$ est subordonné à un objet
galoisien à opérateurs $(P, G)$ de $R_{S}$, de sorte que
\[
  R_{S} \;=\; \varinjlim_{P} R^{P}_{S}
\]
est \emph{limite inductive filtrante} de sous-catégories, et l'on obtient une
description de $R_{S}$ en termes d'un \emph{système projectif} d'objets à
opérateurs $(Z_{i}, G_{i})$.

\begin{quote}
\emph{C'est l'exemple 2 du tapuscrit de 1967}, page 16 du carton : un système
projectif $(Z_{i}, G_{i})$ de schémas à opérateurs, et
$\mathrm{Rev}^{\mathcal{R}}(S) = \varinjlim_{i}
\mathrm{Rev}^{\mathcal{R}_{i}}(S)$. Là encore, l'exemple est devenu une
description.
\end{quote}

Si l'objet final de $R_{S}$ est \emph{connexe}, on peut prendre les $P_{i}$
connexes ; les morphismes de transition $G_{j} \to G_{i}$ sont alors
\emph{surjectifs} et $Z_{i} \simeq Z_{j}/H_{ij}$, et les $Z_{i}/G_{i}$ sont
connexes. Ce sont exactement les conditions sous lesquelles, en 1967, les
foncteurs de transition de l'exemple 2 étaient pleinement fidèles.

\subsection*{Et pourquoi il faut recoller}

Vient alors l'observation qui justifie tout ce qui suit :

\begin{quote}
Partant de $R$ et construisant les $(Z_{i}, G_{i})$, on obtient un nouveau
$(R', r')$ et une inclusion $R' \hookrightarrow R$ ; il n'est \emph{pas} vrai
en général que ce soit une équivalence : \emph{le champ $R$ n'est pas
nécessairement engendré par ses sections globales}. Donc pour définir un $R$ en
général, il faut localiser et \emph{recoller}.
\end{quote}

\begin{quote}
\emph{C'est la raison d'être de l'exemple 3 de 1967}, et c'est ici seulement
qu'elle est dite. La page 35 du carton l'expliquait déjà en termes concrets ---
« sometimes the ramification models are not defined globally (for instance in
the Kummer case, when we give divisors which are not principal), only locally »
--- mais sans énoncé. L'exemple 3, avec ses quatre conditions, puis ses cinq
conditions a) à e) de la page 32, était l'appareil de recollement écrit à la
main. Ici, la nécessité du recollement est un fait sur le champ $R$, et le
recollement lui-même va être fait une fois pour toutes.
\end{quote}

\subsection*{Les homomorphismes}

Un \emph{homomorphisme} de domaines de ramification $(R, r) \to (R', r')$ est
un couple $(\varphi, \lambda)$ où $\varphi : R \to R'$ est un foncteur
\emph{continu et exact} et $\lambda$ un isomorphisme de foncteurs continus
$r'\varphi \simeq r$.

\pagerange{197}{198}
\section*{La 2-catégorie $\mathrm{Domram}(S)$, et le recollement (pages 197 et 198)}

\subsection*{Une 2-catégorie}

Les domaines de ramification sur $S$ forment une \emph{2-catégorie},
$\mathrm{Domram}(S)$ ; et pour $S$ variable, ces 2-catégories forment une
\emph{2-catégorie fibrée}.

C'est ici que se règle ce devant quoi la page 27 avait reculé : « On recule
devant la tâche d'expliciter une propriété de transitivité pour un composé
$U \to T \to S$, et jetant provisoirement un voile pudique sur ces abîmes ».
Les abîmes étaient la 2-catégorialité, et elle est ici assumée comme une
structure plutôt que subie comme une complication.

\subsection*{Proposition 6}

Si $(R, r)$ est défini par $(Z, G)$, on dispose dans $R_{S}$ d'un objet
\emph{galoisien canonique} $P$ de groupe $G$. Un homomorphisme
$(R, r) \to (R', r')$ définit alors un objet $\varphi(P)$ de $R'$, galoisien de
groupe $G$, et un isomorphisme compatible à $G$
\[
  r'(\varphi(P)) \;\overset{\alpha}{\simeq}\; Z .
\]

\textbf{Proposition 6.} Cela donne une \emph{équivalence} entre
$\underline{\mathrm{Hom}}\bigl((R,r),(R',r')\bigr)$ et la catégorie des couples
$(P', \alpha)$ où $P'$ est un objet galoisien de $R'_{S}$ de groupe $G$ et
$\alpha$ un $G$-isomorphisme $r'(P') \simeq Z$.

\textbf{Corollaire.} Dans tout homomorphisme $(\varphi, \lambda)$ de domaines
de ramification, $\varphi$ est \emph{nécessairement pleinement fidèle}.

C'est un énoncé fort, et il faut le mesurer : \emph{il n'y a pas de morphisme
non trivial qui écrase}. Un domaine de ramification ne s'envoie dans un autre
que comme sous-domaine, et le corollaire suivant dit lequel.

\textbf{Corollaire.} Si $(R', r')$ est défini par $(Z', G')$, alors
$\underline{\mathrm{Hom}}\bigl((R,r),(R',r')\bigr)$ est équivalent à la
catégorie des revêtements \emph{principaux} $P'$ de $Z'$ de groupe $G$, sur
lesquels $G'$ opère compatiblement avec ses opérations sur $Z'$, munis d'un
isomorphisme de $G$-schémas sur $S$
\[
  P'/G' \;\simeq\; Z .
\]

\subsection*{Le recollement}

Soient des $S'_{i}$ couvrant $S$, sur chacun un schéma à groupe d'opérateurs
admissible $(Z_{i}, G_{i})$ ; sur chaque $S_{i} \times_{S} S_{j}$ une
\emph{subordination} de $(Z_{ij}, G_{i})$ à $(Z_{ji}, G_{j})$, où
$Z_{ij} = Z_{i} \times_{S_{i}} (S_{i} \times_{S} S_{j})$ ; sur les triples, des
isomorphismes de transitivité ; et enfin une condition de compatibilité sur les
triples. Alors on peut \emph{recoller} les $(R_{i}, r_{i})$.

\subsection*{La construction générale, par les gerbes}

Soit $C$ une catégorie \emph{fibrée} de « modèles » sur $\mathrm{Et}(S)$ ---
pas nécessairement un champ --- et $s : C \to \mathrm{Fl}_{S}$ un foncteur
continu. On suppose :

\begin{itemize}
  \item $C$ fibrée en \emph{groupoïdes} ;
  \item deux objets d'un $C_{S'}$ localement isomorphes ;
  \item $C_{S'}$ localement non vide.
\end{itemize}

Le champ associé est alors une \emph{gerbe}. On suppose de plus que le
\emph{lien} de cette gerbe est \emph{localement constant fini} : pour tout
$x \in C_{S'}$, le faisceau associé au préfaisceau $S'' \mapsto
\operatorname{Aut}(x_{S''})$ est représentable par un revêtement étale de
$S'$.

On a alors une « représentation » de cette gerbe dans le préchamp
$\mathrm{Fl}_{S}$, et l'on associe au tout un domaine de ramification :

\begin{tikzcd}[column sep=small, row sep=small]
& \mathrm{Rev}\,\mathrm{Fl}_S \arrow[d] \\
\mathcal{G} \arrow[r] \arrow[d] & \mathrm{Fl}_S \\
\mathrm{Et}(S) &
\end{tikzcd}

\emph{les objets de $R_{S'}$ sont les homomorphismes continus
$\mathcal{G} \to \mathrm{Rev}\,\mathrm{Fl}_{S}$ qui relèvent
$\mathcal{G} \to \mathrm{Fl}_{S}$.}

\begin{quote}
C'est, littéralement, la définition de 1967 : le lot 1 posait une catégorie
fibrée $\mathcal{C}$ sur $\mathrm{Et}(S)$, un foncteur cartésien $s$ vers les
schémas relatifs, et définissait $\mathrm{Rev}^{\mathcal{R}}(S)$ comme la
catégorie des \emph{sections cartésiennes} de l'image inverse de la catégorie
fibrée des revêtements étales. Ce sont les mêmes objets, et le mot
« relèvement » est le mot juste pour « section cartésienne ». Ce qui a changé
tient en deux mots : $\mathcal{C}$ est maintenant une \emph{gerbe} à lien
localement constant, et le résultat est un \emph{champ}. La normalisation
demandée page 24 est faite.
\end{quote}

\pagerange{199}{200}
\section*{Les exemples kummériens, et la classe une quatrième fois (pages 199 et 200)}

\subsection*{Exemple 1 : des équations globales}

Donnons-nous une famille $\underline{a} = (a_{i})_{i \in I}$ de sections de
$\mathcal{O}_{S}$ et une famille $\underline{n} = (n_{i})_{i \in I}$ d'entiers
$\geqslant 1$. On prend pour $C_{S'}$ le groupoïde correspondant, avec
\[
  s(C_{S'}) \;=\; Z^{\underline{a}}_{\underline{n}} , \qquad
  \mu_{\underline{n}}(S') \;=\; \prod_{i \in I} \mu_{n_{i}}(S') ,
\]
et l'opération qui en dérive. Les axiomes sont vérifiés. Ici la gerbe est
\emph{triviale} --- il y a une section globale, à savoir $\underline{a}$
elle-même --- et le domaine se réduit aux revêtements étales de
$Z^{\underline{a}}_{\underline{n}}$ sur lesquels $\mu_{\underline{n}}$ opère
compatiblement.

\textbf{Exemple 1 bis.} La limite sur $\underline{n}$ : le \emph{domaine de
ramification kummérien} défini par $\underline{a}$.

\subsection*{Exemple 2 : des diviseurs, et la gerbe}

Donnons-nous une famille $(D_{i})_{i \in I}$ de \emph{diviseurs} et une famille
$(n_{i})$ d'entiers. On prend
\[
  \operatorname{Ob}(C_{S'}) \;=\; \bigl\{\, \underline{a} = (a_{i})_{i \in I}
  \;:\; a_{i} \text{ est une équation de } D_{i} \,\bigr\} ,
\]
\[
  \operatorname{Hom}(\underline{a}, \underline{b}) \;=\;
  \bigl\{\, \xi = (\xi_{i})_{i \in I} \;:\; b_{i} = a_{i}\,\xi_{i}^{\,n_{i}}
  \,\bigr\} ,
\]
la composition et les images inverses étant évidentes. Le groupe des
automorphismes d'un objet est $\{\xi : \xi_{i}^{n_{i}} = 1\} =
\mu_{\underline{n}}$, et le champ associé est donc une \emph{gerbe abélienne de
groupe $\mu_{\underline{n}}$}. Elle est localement non vide --- les équations
existent localement --- et sa classe obstrue l'existence d'une section globale.
Grothendieck la nomme :

\begin{quote}
L'existence d'une section dépend de la \emph{nullité} d'un élément de
$H^{2}(X, \mu_{\underline{n}})$, qui n'est autre, comme on devine, que la
\emph{classe de cohomologie des $D_{i}$} dans les $H^{1}(X, \mu_{n_{i}})$ !
\end{quote}

C'est le théorème du lot 5, énoncé pour la quatrième fois dans ce carton, et
sous sa forme la plus nue.\footnote{Les trois autres : le lemme 2 de la page 89
(la classe de l'extension des groupes fondamentaux est celle de
l'auto-intersection de la fibre spéciale) ; le programme de la page 100 (pour
une famille quelconque de diviseurs de Cartier, $\xi_{i} =
c(\mathcal{O}_{X}(D_{i})) \bmod n_{i}$) ; et le corollaire de la page 102. Les
deux exposants de $H$ sur cette page sont tracés de la même main, qui ne
distingue pas le 1 du chiffre romain ; le second est un 2 : la classe d'un
fibré en droites est dans $H^{1}(\mathbb{G}_{m})$, et sa classe de Chern dans
$H^{2}(\mu_{n})$ par la suite de Kummer.}

\textbf{Exemple 2 bis.} La limite sur $\underline{n}$ : le \emph{domaine de
ramification kummérien} défini par $D = (D_{i})$.

\subsection*{Ce que ces exemples sont, aujourd'hui}

La gerbe de l'exemple 2, avec les revêtements étales de son revêtement
kummérien équivariants sous $\mu_{\underline{n}}$, est ce qu'on appelle
aujourd'hui le \emph{champ des racines}
$\sqrt[\underline{n}]{\underline{D}/S}$ : le champ dont les points sont les
choix de racines $n_{i}$-ièmes des équations des $D_{i}$, et dont la gerbe
d'inertie le long de $D_{i}$ est $\mu_{n_{i}}$. Les revêtements de type $R$ ne
sont alors rien d'autre que les revêtements finis étales de ce champ, et le
groupe $\pi_{1}^{\mathcal{R}}(S)$ construit à la page 25 du carton est son
groupe fondamental.\footnote{Les champs des racines sont de Charles Cadman et
d'Abramovich, Graber et Vistoli, dans les années 2000 ; les gerbes sont de Jean
Giraud (1971) ; et l'énoncé que les revêtements modérément ramifiés d'un couple
$(S, D)$ à croisements normaux sont les revêtements étales d'un champ, ou
--- ce qui revient au même --- les revêtements Kummer-étales du schéma
logarithmique associé, appartient à la géométrie logarithmique de Kato, Illusie
et Nakayama, dans les années 1990. Aucun de ces mots n'est disponible ici, et
rien dans cette note n'affirme de filiation : elle constate que l'objet, sa
gerbe, sa classe d'obstruction et la description de ses revêtements sont
construits dans ces pages.}

\subsection*{Le boitement, avoué}

La page se termine sur une réserve : « Ces exemples boitent un tant soit peu ;
on voudrait au fait que les $n_{i}$ ne soient pas premiers avec les
[caractéristiques] résiduelles des $D_{i}$ seulement \dots ». C'est-à-dire :
la condition de modération devrait porter sur chaque $D_{i}$ séparément, et non
globalement sur $S$. C'est exactement l'assouplissement que la donnée
kummérienne du lot 1 opérait en laissant les $u : G_{S'} \to
\mu_{\underline{n},S'}$ varier localement --- et le lot 11 va reprendre la
question sous le titre des \emph{familles génératrices}.

La page s'interrompt en cours de phrase, sur « Dans le cas de l'exemple 2 ». La
suite est au lot 11, et elle ne revient pas sur les exemples.

\end{document}
